\section{Scope Statements}\label{appendix-assumptions}

This appendix lists the principal scope statements for the finite decision-thermodynamic framework:

\begin{itemize}
	\item \textbf{Canonical encoding:} The main theorems are exact for the canonical binary decision problem attached to the bounded decision system, and the canonical object is characterized by an initiality theorem in the exact finite-resolution category over the canonical state space. A noncanonical transport theorem is now proved for optimizer-preserving encoding equivalences that also preserve/reflect coordinate-agreement relations; the same transport layer now has a continuous-state measurable-encoding bridge package (with partition-code compatibility) that transports coordinate relevance, structural rank, and Landauer floors, and exports those transports as one bundled theorem endpoint. An explicit extension-interface witness bundle now captures richer physical encoding transport obligations as theorem-level witnesses.
	\item \textbf{Kernel-completion categorical endpoint:} The proved endpoint theorem is for operation-kernel schemas (operation data plus kernel congruence), with object-level equivalence to bare operation systems, hom-level full-faithful embedding, and universal no-collapse canonicity/factorization package at that level. The AE-germ quotient and an operation-preserving renormalization-flow quotient-invariance law are now included as explicit instances.
	\item \textbf{Measure-kernel dynamics layer:} The theorem package now includes abstract semigroup transport for detailed-balance/stationarity claims, kernel-power and quotient-calculus transport laws, path-space Crooks/Jarzynski lifts, explicit path-measure and process-level transport endpoints, canonical finite-horizon recursion, projective/Kolmogorov interfaces, a canonical truncation-witness instantiation, and concrete Giry-bind measurable-kernel instantiation results. The canonical finite-horizon truncation measurability step is now proved constructively in Lean (no truncation axiom). A continuous-time/continuous-state interface endpoint is also included at theorem level (continuous-time kernel semigroup discretization to additive scale flows, plus measurable two-time path projections). The finite-horizon constructive bridge additionally exports direct projective-marginal and Kolmogorov-extension marginal recovery formulas, and the first-principles layer now includes both a microscopic-derivation constructor and a concrete molecular-Hamiltonian+thermostat constructor that map microscopic drift inequalities directly into the first-principles Langevin interface, together with a direct microscopic dissipativity-to-ergodicity export theorem, a canonical continuous path-process closure theorem that packages process regularity/uniqueness certificates with constructive finite-horizon closure data, and an explicit joint microscopic+canonical-path closure bundle theorem. The newest closure layer now also provides a constructive stochastic-analysis discharge bundle that instantiates canonical process regularity directly from finite-horizon bridge data.
	\item \textbf{Structural chain:} Counting Gap fixes the finite-event statement. Bounded Acquisition fixes the traversal-rate statement. Discrete Acquisition and One Transition, One Bit fix the acquisition-event interface.
	\item \textbf{Landauer calibration:} Thermodynamic cost is calibrated by a per-bit Landauer floor. Stronger substrate-dependent lower bounds may exist, but they are additional assumptions, not part of the theorem package under discussion.
	\item \textbf{Quantum substrate instantiation layer:} The abstract quantum readout/decoherence theorem now has a concrete spin-$\tfrac12$ instantiation endpoint: explicit two-state readout witness, exact canonical rank-one embedding, and exact $k_B T\ln 2$ decoherence event identity. A theorem-level extension witness bundle now includes open-system channel modeling obligations explicitly, reducing richer channel extensions to witness discharge.
	\item \textbf{Exact-resolution setting:} The results concern exact sufficiency and exact-resolution cost. Approximate, stochastic, or bounded-confidence regimes require separate analysis.
	\item \textbf{Composed force-field layer:} The theorem package now includes an explicit composed classical force-field interface (bonded + Lennard-Jones + Coulomb) with architecture/decision-problem projections, shellwise Lipschitz composition, and concrete biomolecular calibration bundles for zero-shell, anchor-nontrivial, calibrated full-state interaction, and pairwise-geometric families with explicit-geometry-derived shell constants (all with half-gap transport endpoints). A realism-augmentation layer now additionally supports explicit solvent, polarization, many-body, and long-range correction terms with summed-shell Lipschitz transport and half-gap endpoint export, and a constructive empirical instantiation layer now anchors these shells to dataset-level mean/std-error calibrations with explicit $z/\sqrt{N}$ finite-sample shell envelopes. The same layer now also includes a concrete reference-realism calibration package with explicit shell constants, explicit aggregate-shell value export, and direct constructive transport theorems, plus an electronic-structure correction layer (charge transfer + metal coordination) with explicit shell/error envelopes and corresponding Lipschitz/error/half-gap transport endpoints, including a reference specialization where the electronic shell vanishes and half-gap transport reduces to the realism-only shell, a QM-grounded chemistry-class envelope layer providing tighter shell/error constants for designated active chemistry classes, and a QM-tight-shell half-gap transport endpoint. A higher-order multipole closure theorem now exports explicit multipole shell-Lipschitz constants under constructive witness discharge.
	\item \textbf{Finite state family:} The entropy and replication theorems are finite counting results. Continuum models must first be reduced to a finite decision quotient before these arguments apply.
	\item \textbf{Fisher observation model:} The canonical explicit optimal-action observation channel remains the fully constructed model, and the same section now includes a theorem-level general observation-channel Fisher interface endpoint: any channel with declared coordinatewise Fisher-to-relevance identification inherits the structural-rank Fisher sum identity and diagonal relevance recovery. The extension-interface witness bundle now packages noisy/partial observation-channel obligations explicitly, keeping concrete experimental-channel instantiation at witness-discharge level.
	\item \textbf{Admissibility canonicality layer:} The fixed-erasure-count uniqueness theorem is stated in the rank-normal-form collapse category (objects are collapsed ranks, morphisms are rank-monotone maps). Lifting this uniqueness to richer state- or action-level isomorphism notions requires additional structure.
	\item \textbf{Collapsed-rank categorical layer:} The proved category structure is the preorder/rank-normal-form layer (identity/composition, initial object, min/max product-coproduct, additive tensor, monotone rank functor). In the unbounded setting, no terminal object exists; in the bounded-rank slice, terminal objects, finite limits/colimits, finite-family meet/join objects, arbitrary-family meet/join objects (complete-lattice form), and monotonicity/idempotence/nonempty-absorption corollaries are recovered via the bounded index-order encoding, together with binary meet/join rewrite identities (commutativity, associativity, idempotence, absorption). No topos-level claim (e.g., subobject classifier) is asserted here.
	\item \textbf{Combinatorial lower-bound layer:} The binary optimizer-richness floor remains fully proved, and a nonbinary/VC transfer endpoint is now included: under a declared $q$-ary class envelope and growth witness, structural rank inherits the same lower bound. Concrete nonbinary model instantiations therefore reduce to proving the declared $q$-ary class envelope.
	\item \textbf{Proofreading specificity layer:} The Hopfield--Ninio reduction is stated for the explicit exponential specificity model and a declared overhead-energy witness under Landauer calibration, with a kinetic branch-rate specialization requiring positive equilibrium correct/error rates and an explicit branch-ratio calibration identity; richer kinetic-network assumptions must still be bridged into that witness interface.
	\item \textbf{Fluctuation-relation layer:} The Crooks/Jarzynski reduction is stated for finite strictly positive trajectory distributions with explicit calibration identities, includes the detailed-balance equilibrium calibration corollary, and includes an explicit nonequilibrium molecular-dynamics class interface theorem that derives Crooks standard form from declared stepwise rank calibration plus cumulative work calibration. A continuous overdamped-Langevin-to-discrete-MCMC bridge endpoint is now added at theorem level (Boltzmann detailed balance plus Euler--Maruyama discretization error witness), together with finite-state measure-theoretic Boltzmann stationarity derivation from detailed balance + row-stochasticity, continuous-state measure-kernel closure, an explicit Ito/Fokker--Planck/Harris bridge layer, and an infinite-dimensional path-measure layer that now has both certificate-export and constructive finite-horizon/Kolmogorov-extension derivation endpoints.
	\item \textbf{Physics-assumption interfaces:} Thermodynamic/kinetic premises used by these bridges are represented as explicit proposition interfaces carried as theorem hypotheses, and the current build now includes concrete witness/proof instantiations for multiple interfaces (EP1/EP2/EP3 locality witnesses, reversible-chain entropy-production nonnegativity witness, FDT-to-stationarity discharge theorem, differentiability-to-Born--Oppenheimer discharge theorem, and large-radius Lennard-Jones lattice-tail convergence witnesses), plus bundled paper4-to-paper3 witness-import theorems. The discharge layer now also exports direct certificate-to-interface conversion for TUR inequalities and velocity-Verlet shadow-Hamiltonian cubic-drift bounds, resolves the exploratory stochastic-preservation relevance conjecture on Boolean product spaces under full-support hypotheses, and adds unrestricted-distribution closure theorems under support-transport, primitive finite-time dynamics, and explicit one-step dynamics derivations, including an explicit-step-to-support-transport assumption-reduction path. The active paper3/paper4 proof sources no longer rely on custom global Lean \texttt{axiom} declarations for these premises.
	\item \textbf{Speed-accuracy/on-rate layer:} The admissibility-indexed speed law is stated for exact-resolution witnesses in bounded regions with positive signal speed. The on-rate envelope additionally assumes positive time horizon and positive relaxed structural rank so the reciprocal rate expression is well-defined. A protocol-noise theorem layer now adds explicit finite-sample absolute-error budgets and margin conditions under which threshold/ranking inferences for kinetic observables are guaranteed stable, plus concentration/identifiability wrapper layers that export those same conclusions together with statistical certificate predicates; a replicate-aware identifiability extension further incorporates systematic-bias-aware pathway margins and replicate-count certificates, and a hierarchical multi-dataset extension adds pooled-replicate finite-sample-rate metadata with shared-bias margin transport and explicit chemical/kinetic rate-margin theorem endpoints, including dataset-uniform rate-margin transport theorems, external rate-constant upper-bound reduction theorems, required sample/replicate complexity-inversion transport theorems, and a joint chemical+kinetic required-size bundle theorem. The same section now also includes a unified thermodynamic/kinetic single-model theorem bundle, a universality/OOD uncertainty-calibration layer, and a unified physical/OOD/prospective empirical closure bundle theorem. The zero-collapse speed specialization requires the additional backward factorization witness at the same tolerance increment.
	\item \textbf{Chemical-state and empirical closure layer:} Chemical microstate effects are now represented both as static coupled utility contributions and as unified finite-state dynamic layers carrying joint protonation/tautomer/ionic/solvent/water-bridge metadata with pH/ionic window certificates and stationary expected-utility transport. The same build also includes an absolute binding free-energy correction stack (standard-state, finite-size, long-range, net-charge, restraint) with summed correction-bound transport and a prospective blinded-benchmark empirical-closure layer that exports calibration/predictive coverage and failure-rate-bound certificates.
	\item \textbf{Allosteric distance-decay layer:} The shell-sum theorem is finite and profile-driven; the exponential specialization assumes nonnegative envelope scale and decay rate and a shell-cardinality witness bounded by the ceiling of that exponential envelope, while the polynomial specialization assumes nonnegative envelope scale, natural exponent profile, and the corresponding ceiling-bounded shell witness. The polynomial series-budget corollary additionally requires an explicit finite bound on the partial reciprocal-power sum used in the closed-form real inequality.
	\item \textbf{Replication theorem:} The finite replication entropy gap is a theorem of the calibrated exact-resolution model. England's 2013 theorem belongs to a stochastic-thermodynamic path-space model.
	\item \textbf{Finite-budget theorem:} Finite-Budget No-Collapse is a theorem about bounded budgets, positive per-bit cost, and exponential lower-bound growth.
	\item \textbf{Top-level executability layer:} The development now includes a constructive rational-certificate top-level path from molecular input to ArrayDSL JSON export with theorem-level computability statement, plus a normalization/deprecation bridge theorem family. Automatic wrapper theorems now derive legacy/constructive alignment witnesses directly from constructive outputs, yield a fully constructive deprecation-readiness theorem for the top-level pipeline, provide a proposition-level equivalence principle for constructive-only normalized outputs, and include direct migration lemmas for both core and extended downstream consumers (membership/cardinality/refinement-payload/export predicates). The legacy noncomputable path remains available only for compatibility with older statement forms.
	\item \textbf{Latest physical-completeness discharge layer:} The current build now also includes theorem-level closures for (i) concrete Hamiltonian+bath to end-to-end Langevin endpoints, (ii) canonical path-process construction directly from finite-horizon bridge data, (iii) method-specific affine QM calibration transport to tight electronic shell/error bounds, (iv) conditioned-mechanism-derived unified chemical dynamics with stationary fixed-point identity, (v) protocol-derived absolute free-energy correction closure, (vi) simulator-derived thermo/kinetic bundling from shared model measurements, (vii) transfer-calibration-derived OOD uncertainty bounds, and (viii) count-order-derived finite-sample inversion monotonicity discharge for hierarchical chemical/kinetic required-size theorems. The newest extension layer further adds protocol-derived active-class QM calibration transport, a Gibbs-factor-derived chemical mechanism closure with explicit stationary law formula, square-count-complexity-derived finite-sample inversion transport, explicit Hamiltonian-elimination microscopic drift export, explicit constant-drift molecular SDE endpoint constants, constructive generator-residual-to-canonical-path closure, workflow-specific QM decomposition transport, reversible detailed-balance chemical dynamics bundling, trajectory-estimator absolute free-energy correction closure, quantified simulator-control thermo/kinetic bundling, mechanistic OOD envelope transfer bounds, and model-dependent finite-sample inversion constants with joint hierarchical required-size discharge. The latest pass additionally upgrades these with finite-difference Hamiltonian drift derivation, non-constant realistic molecular SDE closure, coefficient-derived canonical regularity construction, barrier-crossing Kramers/Eyring reversible kinetics, descriptor-calibrated mechanistic OOD transfer with prospective protocol witness, minimax-optimality finite-sample transport, and a bundled constructive stochastic-analysis plus multipole scope-gap discharge theorem. The current extension now further contributes explicit Ito-Wiener-filtration semantics, Hamiltonian $h\to0$ Mori--Zwanzig limit closure, force-field-derived SDE constants, benchmark-derived QM workflow transport, potential-landscape prefactor-derived barrier kinetics, spectral-gap concentration-to-trajectory correction transport, integrator-stack-derived simulator controls, learned-descriptor prospective OOD generalization transfer, estimator-backed minimax derivation closure, and explicit scope-extension witness bundling.
\item \textbf{External validation evidence layer:} The theorem package now includes an explicit three-way external-validation integration layer: (i) pre-registered prospective benchmark superiority over declared strong baselines on calibration/ranking/kinetics-free-energy consistency, (ii) independent outside-team replication with matched protocol and reproduced superiority clauses, and (iii) strict downstream campaign improvements (hit identification, triage quality, campaign efficiency). A formal core-gap dismissal predicate is defined, and the integrated bundle theorem proves its negation once the primary (1)+(2) evidence core is instantiated. The latest pass additionally introduces fixed-contract benchmark locking (metrics/datasets/splits/blind rule/baseline roster), signed artifact provenance records, separate-team/separate-compute independent-replication provenance closure, randomized/blinded/balance-controlled downstream causal-quality closure, and one concrete artifact instantiation theorem for the full three-way package.
\item \textbf{Pass-29 measure-theoretic derivation layer:} The current pass adds a measure-theoretic Langevin endpoint bridge (probability path-law semantics with measure-level stationarity/ergodicity transport), a constructive Ito/Wiener derivation layer where adaptedness/integrability/Ito/martingale clauses are extracted from explicit objects, an explicit-profile PDE/operator derivation layer feeding canonical regularity closure, a microscopic open-system/noisy-observation scope-derivation layer, and a numerical-stack theorem identifying top-level simulator control flags with explicit mismatch-vs-tolerance inequalities.
\item \textbf{Pass-29 attested external-evidence layer:} External validation is further strengthened with verifier-checked signatures, immutable-manifest ingestion records, timestamp-based preregistration locking, digest-based contract/result binding, typed team/compute provenance backed by signed identity/compute/protocol/execution artifacts, and measured randomized/blinded/balance diagnostics for downstream causal isolation, while preserving the no-credible-dismissal core-gap conclusion at concrete attested instantiation level.
\item \textbf{Pass-30 immutable-store ingestion strengthening:} The attested external-evidence stack is now additionally tied to an explicit immutable artifact store model (URI fetch + digest function), with theorem-level witnesses that each concrete attested artifact manifest is store-backed by payload/digest consistency and that the full attested three-way external-validation no-credible-dismissal conclusion persists under this store-backed ingestion layer.
\item \textbf{Pass-31 reviewer-response closure layer:} The theorem package now explicitly states that exact docking decision utility is inherited directly from the molecular binding problem object, identifies docking structural rank with native-coordinate relevance count, proves a physical $3(N-k)$ rank budget under cutoff-locality-relevant atom counts, proves binary summary rank monotonicity relative to full pose-selection docking, adds a resolver-free equilibrium bridge combining docking free-energy floor with detailed-balance path-ratio unity, and adds falsifiable independent-rank prediction theorems linking independently certified rank lower bounds to both equilibrium affinity upper bounds and necessary one-hop contact-shell geometry budgets.
\item \textbf{Pass-32 physical-completeness closure layer:} The theorem package now additionally includes (i) per-system concrete-physics witness discharge bundles that derive strict-optimality, cutoff-bounded perturbation behavior, and both excluded-atom/contact-shell docking rank budgets from explicit force-field/cutoff/geometry certificates; (ii) an end-to-end measurement-free equilibrium chain from partition-ratio free-energy definitions plus explicit standard-state/finite-size/long-range/net-charge/restraint correction budgets to the exact $\Delta G_{\mathrm{drive}}$ witness used in rank-based affinity bounds; (iii) independent structural-rank interval extraction theorems from relevance/sufficiency certificates with no affinity fit input; (iv) pre-registered hard-threshold falsification protocol logic with explicit fail conditions; (v) finite-sample/high-confidence margin-to-true-violation transport theorems for assay-noise handling; and (vi) chemistry-condition data-backed uncertainty propagation endpoints from calibration windows/posterior constants through free-energy correction error budgets to final equilibrium $K_d$ intervals.
\item \textbf{Pass-33 autonomous execution closure layer:} The theorem package now adds an execution-level closure stack that formalizes: (i) per-case real Hamiltonian/force-field/geometry artifact discharge into concrete docking witnesses for all declared targets, (ii) certified numerical partition-function and correction-term closure feeding directly into equilibrium $K_d$ bounds, (iii) production independent-srank extractor outputs with explicit structure/mechanics-only and no-affinity-input certificates, (iv) locked prospective falsification-run records with split/metric/blind locking and fail-condition-to-falsification soundness, (v) assay-noise calibration models decomposed by replicate/batch/instrument error with both fail-call and pass-call high-confidence validity transport, (vi) target-class chemistry completeness for protonation/tautomer/ionic/solvent/water/electronic calibration uncertainty propagated to final $K_d$ intervals, and (vii) outside-team/outside-compute replication-at-scale packaging that exports the full target-level closure proposition pointwise across new targets.
\item \textbf{Pass-34 instantiation-constructor closure layer:} The theorem package now additionally provides operational constructor endpoints that make the pass-33 objects directly instantiable in canonical low-assumption form: (i) one-state exact partition constructors with certified positivity, (ii) canonical zero-correction one-sample calibration closure, (iii) upper-only independent-srank extractor templates from supplied sufficiency certificates (empty lower set + explicit no-affinity provenance), (iv) canonical zero-rank single-target full-closure instance construction from real witness package plus upper sufficiency certificate, and (v) attested-provenance external-replication constructors including a singleton-target concrete attested replication bundle.
\item \textbf{Pass-35 computable-pose and RMSD-probability closure layer:} The theorem package now additionally formalizes (i) an explicit finite-enumeration pose-search algorithm endpoint (list-argmax over covered action families) with exact optimizer-membership guarantee, (ii) finite RMSD-success probability semantics from posterior mass over threshold-satisfying poses with unit-interval calibration bounds, (iii) monotone transport from paper4 finite top-$k$ posterior mass to RMSD-success probability under RMSD-coverage certificates, and (iv) a bundled RMSD-probability-derived pose-solver endpoint combining selected-pose threshold validity with certified lower bounds on success probability.
\item \textbf{Pass-36 definitive raw-endpoint closure layer:} The theorem package now additionally closes an end-to-end raw pocket+ligand sampled-docking endpoint in Lean by formalizing: (i) canonical raw-input sampled-family construction (grid-lifted candidate family plus base-action anchor), (ii) posterior synthesis and normalization certificates from that raw constructor, (iii) benchmark-vs-deployment contract split with conservative proxy-RMSD calibration and deployment-to-benchmark implication, (iv) canonical ProgramIR execution witnesses refining runtime accept/reject flags to accepted/failure solver certificates, and (v) canonical raw benchmark/deployment solver wrappers with acceptance equivalences, totality theorems, and a bundled definitive endpoint theorem.
\item \textbf{Pass-37 definitive API consolidation layer:} The theorem package now further consolidates the raw endpoint into single-name definitive API wrappers with theorem-level guarantees: (i) canonical runtime-output interpreter refinement to benchmark accepted/failure certificates without external execution assumptions, (ii) definitive deployment acceptance equivalence both to benchmark and canonical deployment contracts, (iii) definitive benchmark/deployment totality, and (iv) a final full-closure bundle theorem for the top-level definitive raw cross-docking API.
\item \textbf{Pass-38 acceptance-flag exactness layer:} The theorem package now additionally proves exact two-sided accept/reject characterizations for the definitive API: benchmark accepted/failure iff benchmark-contract/negation, deployment accepted/rejected iff deployment-contract/negation, runtime accept-flag equivalence to both benchmark acceptance and deployment acceptance, runtime reject-flag equivalence to deployment rejection, and consolidated top-level closure bundling these exactness properties.
\item \textbf{Pass-39 computable rational-kernel layer:} The theorem package now additionally provides a computable (Boolean-decidable) rationalized acceptance kernel with explicit margin inequalities, together with soundness transport from rational checks to real benchmark contract satisfaction and an acceptance-refinement theorem from computable flag truth to formal accepted solver certificate existence.
\item \textbf{Pass-40 interpreter/report refinement layer:} The theorem package now additionally formalizes an explicit canonical interpreter-state runtime view and proves its refinement equivalence to solver certificate outcomes, plus definitive wrapper theorems for interpreter-runtime identity and consolidated report-field runtime-accept/deployment-accept equivalence.
\item \textbf{Pass-41 complete Lean bundle layer:} The theorem package now additionally exposes one top-level complete Lean closure theorem for the definitive raw cross-docking API, jointly bundling support inclusion, JAX codegen success, benchmark/deployment contract equivalences, runtime accept/reject exactness, deployment totality, and computable rational-kernel acceptance refinement to benchmark certificates.
\item \textbf{Pass-42 report reject-flag exactness layer:} The theorem package now additionally proves report-level runtime reject-flag equivalence to deployment rejection certificates, complementing the existing report-level runtime accept-flag equivalence and completing two-sided report runtime/deployment exactness.
\item \textbf{Pass-43 endpoint-stabilization layer:} The theorem package now additionally restores full Lean compilability of the definitive raw cross-docking endpoint after generic-ProgramIR integration edits by (i) proving the generic two-fuel entry execution theorem via explicit runtime-flag case analysis and (ii) re-establishing exact-rational witness flag/benchmark-contract equivalence with a direct decide-and-cast proof path; no new paper handles are introduced in this stabilization pass.
\item \textbf{Pass-44 constructive endpoint closure layer:} The theorem package now additionally introduces artifact-instantiated constructive top-level benchmark/deployment decision paths for the definitive raw endpoint (no noncomputable marker on the new top-level constructive wrappers), proves accepted-decision refinement back to legacy benchmark/deployment acceptance certificates, and tightens rational witness assumptions into explicit artifact payload records (artifact identity, backend name, generated-op digest, rational values, and certified error bounds), together with exact-rational artifact equivalence/refinement theorems.
\item \textbf{Pass-45 legacy-chain cleanup layer:} The theorem package now additionally isolates all prior nonconstructive definitive endpoint operators behind explicit `legacy...` names (solver outputs, runtime flag, interpreter/runtime report constructors), keeps theorem compatibility by private in-file aliases only, and exposes public constructive-first decision endpoints plus direct alias-exactness and decision-to-legacy refinement theorems; this removes legacy operators from the public endpoint surface while preserving backward theorem continuity.
\item \textbf{Pass-46 constructive certification and signed-artifact closure layer:} The theorem package now additionally upgrades the public constructive endpoint from decision-only outputs to certificate-carrying benchmark/deployment result types (accepted/rejected certificate objects with exact accepted/rejected iff decision theorems and totality), adds signed-manifest/signature-verifier artifact interfaces with direct conversion into constructive kernel instantiations and acceptance-refinement theorems, and extends exact-rational closure with constructive rejection refinement to legacy benchmark-failure and deployment-rejection certificates (including signed exact-rational transport).
\item \textbf{Pass-47 byte-verified cryptographic closure and constructive-surface rebasing layer:} The theorem package now additionally formalizes an explicit byte-level signed-artifact envelope parser (length-prefix encoding with parse-roundtrip theorem), instantiates a concrete checksum signature verifier and proves end-to-end parse+verify success, lifts that concrete verifier to signed rationalized artifacts, rebases definitive theorem-facing names/claims to certificate-backend terminology (removing explicit legacy naming from the constructive theorem surface), and adds strict-separation rejection transport for broad signed rationalized artifacts, yielding certified benchmark-failure and deployment-rejection refinements beyond exact-rational-only rejection closure.
\end{itemize}

\begin{remark}[Scope Statement: Molecular Independence Hypothesis]\label{rem:molecular-independence-scope}
Corollary~\ref{cor:holonomic-landauer-floor} proves the finite constrained-molecular transport once the RATTLE holonomic topology supplies $k$ independent constraints and the corresponding binary status interface. Corollary~\ref{cor:bond-family-holonomic-floor} derives that independent-count antecedent for concrete finite bond families carrying the explicit certificate $|F|\le N$ with $N>0$. Corollary~\ref{cor:jacobian-holonomic-floor} discharges the Jacobian hypotheses for nonlinear geometric families that provide a pivot-column Jacobian certificate at the chosen configuration. The new geometric-constraint decision-interface theorem additionally covers broad bond/angle/dihedral composites whenever an explicit decision procedure certifies the strict independent-count inequality from the raw constraint data.
\end{remark}
